Methods for annotating spatial transcriptomics data increasingly use a spot's tissue neighbourhood as input, but they are benchmarked within one tissue and condition, so it is unknown whether the spatial signal they learn holds when tissue organization changes. We build a small benchmark from the public Slide-seqV2 testis pucks of Chen et al.\ (2021): three wild-type mouse, three diabetic \emph{ob/ob} mouse, and two human pucks (252,898 beads), with cell-type labels from NMFreg (non-negative matrix factorization regression). Using germ-cell stage annotation in seminiferous tubules as the task, we compare an expression-only classifier with neighbourhood-augmented features, a GraphSAGE model, and test-time smoothing of the expression-only posterior, under held-out-puck, wild-type-to-\emph{ob/ob}, and mouse-to-human transfer. Spatial context adds about one macro-F1 point within a species and up to three on human pucks. The \emph{ob/ob} pucks are not a distribution shift for this task. Across species, expression-only accuracy drops by ten points; neighbourhood-augmented features and test-time smoothing keep their gains, whereas GraphSAGE, which matches them within a species, becomes the worst model, losing five points and fifteen F1 points on round spermatids. The gain from context tracks how informative the neighbourhood is relative to the bead itself rather than the target tissue's neighbourhood homogeneity. Fixed spatial aggregation transfers across the shifts we tested where learned spatial features do not.
